% ============================================================
% 六、个人见解 (Personal Insights)
% ============================================================
\section{个人见解}

\subsection{方法总结}

简要总结本文方法的核心思想与关键步骤（用自己的话概括，5--8 行即可）。

\subsection{优点与适用场景}

\begin{itemize}
    \item \textbf{理论上的优势}：【例如：收敛性有严格保证、假设条件宽松等】
    \item \textbf{实践上的优势}：【例如：实现简单、参数少、调参友好等】
    \item \textbf{特别适用的场景}：【描述方法在何种条件下表现最好】
\end{itemize}

\subsection{局限性与改进方向}

\begin{itemize}
    \item \textbf{理论局限}：【例如：强凸假设在实际中不成立时方法可能退化】
    \item \textbf{实践局限}：【例如：高维情况下计算代价过高、对初始化敏感等】
    \item \textbf{可能的改进方向}：
    \begin{enumerate}
        \item 【自适应步长策略】
        \item 【结合二阶信息的加速方案】
        \item 【推广到更一般的问题类】
    \end{enumerate}
\end{itemize}

\subsection{与其他方法的比较}

简要对比本文方法与文献中的相关方法，突出本质区别（不仅是数值好坏的比较，更要讨论\emph{数学结构上的差异}）。

\begin{table}[htbp]
    \centering
    \caption{与相关方法的比较}
    \label{tab:comparison}
    \begin{tabular}{lccc}
        \toprule
        方法 & 核心数学工具 & 优势 & 局限 \\
        \midrule
        本文方法      & --- & --- & --- \\
        文献 A        & --- & --- & --- \\
        文献 B        & --- & --- & --- \\
        \bottomrule
    \end{tabular}
\end{table}

\subsection{与课程内容的呼应}

反思本文所用数学方法在课程知识体系中的位置：

\begin{itemize}
    \item 本文的数学工具是课程中哪些主题的自然延伸？
    \item 课程中的哪些内容帮助了你理解该问题？
    \item 课程讲述的方法在本问题上有何局限？（是否存在课程方法直接套用失败的情况？你是如何调整的？）
\end{itemize}
